\subsection{Ranked models and rational closure}
\label{sec:p08-ranked-models-rational-closure}

Rational Monotonicity restricts which additional information may defeat a
conditional conclusion. A ranked model represents one rational relation;
rational closure selects a particular extension of a conditional knowledge
base. These are different constructions. The presentation here follows
Lehmann and Magidor's author version \texttt{cs/0202022v1}, especially
\S\S2.6, 3.6, 5.4, 5.6 and 5.8, and Appendix~A. Its 58-page author
pagination is not the 60-page journal pagination
\cite{LehmannMagidor1992,LehmannMagidor2002Author}.

\begin{hbdefinition}{def:p08-rational-language}{The Lehmann--Magidor language and rational relation}
Let \(\mathcal L\) be the finitary classical propositional language over
a set of atoms and \(\mathcal U\) the set of all its Boolean valuations.
Material implication is \(\to\). Classical \(\models\) in this
presentation is validity over all of \(\mathcal U\), unlike a
background-restricted universe in the earlier KLM presentation.
An assertion \(\alpha\mathrel{|\!\sim}\beta\) is a flat pair of
formulas, not a nestable formula connective. A rational relation is a
preferential relation satisfying Rational Monotonicity: if it contains
\(\alpha\mathrel{|\!\sim}\gamma\) but does not contain
\(\alpha\mathrel{|\!\sim}\neg\beta\), it contains
\(\alpha\wedge\beta\mathrel{|\!\sim}\gamma\).
The nonmembership premise is a property of the whole relation, not
negation of an object-language assertion or failure to find a proof in a
finite search. The source uses the six P principles Left Logical
Equivalence, Right Weakening, Reflexivity, And, Or and Cautious
Monotonicity; S and Cut are derived. The earlier native C-plus-Or basis
is equivalent by Theorem~\ref{thm:p08-klm-basis}.
A hard constraint \(\delta\), when needed here, is represented by the
assertion \(\neg\delta\mathrel{|\!\sim}\bot\), not by the
default \(\top\mathrel{|\!\sim}\delta\).
\end{hbdefinition}

\begin{hbdefinition}{def:p08-rational-ranked}{Modular preference and smooth ranked models}
A ranked model \(W=\langle V,l,\prec\rangle\) has a set of states
\(V\), a label map \(l:V\to\mathcal U\), a totally ordered set
\(\Omega\), and a map \(r:V\to\Omega\), with
\(s\prec t\) exactly when \(r(s)<r(t)\).
Labels may have copies. Write \(s\equiv\alpha\) for
\(l(s)\models\alpha\), \(\widehat\alpha\) for these states,
and \(\alpha\mathrel{|\!\sim}_W\beta\) when all minimal
states of \(\widehat\alpha\) satisfy \(\beta\).
Require smoothness on every formula-definable state set. Equivalently,
each nonempty \(r[\widehat\alpha]\) has a least element.
No well-ordering of \(\Omega\), no ordinal-valued ranking, and no
well-foundedness of the entire state order are imposed.
The order is modular when \(x\prec y\) implies, for every \(z\),
\(z\prec y\) or \(x\prec z\).
\end{hbdefinition}

\begin{hbtheorem}{thm:p08-rational-modularity}{proved=annex:p08.rational.modularity}{The four equivalent ranking conditions}
For a strict partial order on a set, the following are equivalent:
incomparable elements have the same strict predecessors; the order is
modular; its complement is transitive; and the order is induced by a
map to a strict total order. Incomparability is then an equivalence
relation and the quotient supplies such a ranking. In a ranked model,
minimal states of a nonempty set are exactly its states of least rank;
formula smoothness is therefore equivalent to the stated least-rank
condition. This supplies the proof omitted for source Lemma~14.
\end{hbtheorem}

\begin{hbtheorem}{thm:p08-rational-soundness}{proved=annex:p08.rational.soundness}{Ranked soundness and the nonmembership witness}
Every smooth ranked model defines a rational relation. For Rational
Monotonicity alone, ranking suffices: its negative premise supplies an
actual minimal \(\alpha\)-state satisfying \(\beta\), and every
minimal \(\alpha\wedge\beta\)-state is then minimal for
\(\alpha\). Smoothness is still needed for the full P soundness
claim. An empty antecedent has every consequent, and cannot supply the
negative premise.
\end{hbtheorem}

\begin{hbdefinition}{def:p08-rational-normal}{Consistent formulas and native exception comparison}
Fix a rational relation. A formula is consistent for it when
\(\alpha\mathrel{|\!\not\sim}\bot\); this is not the same
as classical satisfiability. Let \(S\) be the set of such formulas.
A normal world for \(\alpha\) satisfies every \(\beta\) such
that \(\alpha\mathrel{|\!\sim}\beta\). For \(\alpha,\beta\in S\),
define the source comparison
\[
 \alpha\mathcal R\beta
 \quad\Longleftrightarrow\quad
 \alpha\vee\beta\mathrel{|\!\not\sim}\neg\alpha.
\]
The source's \(\alpha\sim\beta\) in the following quotient means
\(\alpha\mathcal R\beta\) and \(\beta\mathcal R\alpha\),
not the conditional assertion, which has the distinct hooked bar.
Let \(\overline\alpha\) be its equivalence class and \(E\) the
set of classes. Set \(\overline\alpha\leq\overline\beta\)
iff \(\alpha\mathcal R\beta\), and use \(<\) for its strict
part. These class orders differ from the earlier P formula preorder.
\end{hbdefinition}

\begin{hblemma}{thm:p08-rational-transport}{proved=annex:p08.rational.transport}{The comparison construction is well defined}
The relation \(\mathcal R\) on \(S\) is reflexive, transitive and
total. It makes \(E\) a totally ordered set. A formula has a normal
world iff it is consistent for the relation. If
\(\alpha\mathcal R\beta\), a normal world for \(\alpha\)
satisfying \(\beta\) is normal for \(\beta\).
If \(\overline\beta<\overline\alpha\), then
\(\beta\mathrel{|\!\sim}\neg\alpha\).
If a normal world for \(\alpha\) satisfies \(\beta\), then
\(\beta\in S\) and
\(\overline\beta\leq\overline\alpha\).
All required derived P rules and the rational comparison inference are
proved in the annex before the quotient is used.
\end{hblemma}

\begin{hbdefinition}{def:p08-rational-canonical}{The native canonical ranked structure}
Put
\[
 V=\{\langle m,\alpha\rangle\in\mathcal U\times S:
                         m\text{ is normal for }\alpha\},\qquad
 l(\langle m,\alpha\rangle)=m,
\]
and rank \(\langle m,\alpha\rangle\) by
\(r(\langle m,\alpha\rangle)=\overline\alpha\in E\).
Thus \(\langle m,\alpha\rangle\prec\langle n,\beta\rangle\)
iff \(\overline\alpha<\overline\beta\).
Every normal-world/formula pair is included. Normal-world existence,
total ordering of the classes, and smoothness are proved dependencies,
not assumptions inserted into this construction.
\end{hbdefinition}

\begin{hblemma}{thm:p08-rational-canonical-minima}{proved=annex:p08.rational.minima}{Canonical minima and smoothness, including inconsistent formulas}
In the canonical structure a state \(\langle m,\alpha\rangle\)
is minimal in \(\widehat\beta\) iff
\(m\models\beta\) and
\(\overline\alpha=\overline\beta\).
If \(\beta\notin S\), then \(\widehat\beta\) is empty.
For consistent \(\beta\), a nonminimal state above its class has an
explicit minimal predecessor \(\langle n,\beta\rangle\), where
\(n\) is normal for \(\beta\). The structure is therefore smooth
on every definable set.
\end{hblemma}

\begin{hbtheorem}{thm:p08-rational-general-representation}{proved=annex:p08.rational.representation}{General ranked representation in both directions}
A relation on \(\mathcal L\) is rational iff it is the consequence
relation of a smooth ranked model over \(\mathcal U\).
The canonical structure represents it exactly, with a minimal
counterstate for every missing assertion. If the atom set is finite,
a finite ranked representation exists: restrict formula coordinates
to representatives of classical equivalence classes. This finite
quotient step is necessary, since the set of syntactically distinct
formulas remains infinite even over finitely many atoms.
\end{hbtheorem}

\begin{hbtheorem}{thm:p08-rational-nonwellfounded}{proved=annex:p08.rational.nonwellfounded}{Why ranks need not be ordinals}
There is a smooth ranked model on atoms \(p_i\), \(i\in\omega\),
whose rational relation has no well-founded preferential representation.
Take states \(s_i\), \(i\in\omega\cup\{\infty\}\), with
\(s_i\prec s_j\) iff \(i>j\), and
\(l(s_i)(p_j)=1\) iff \(j\geq i\).
The state \(s_\infty\) satisfies none of the atoms and is below
every finite-index state. The relation contains
\(p_i\mathrel{|\!\sim}p_{i+1}\) and
\(p_{i+1}\mathrel{|\!\sim}\neg p_i\) for every \(i\),
but not \(p_0\mathrel{|\!\sim}\bot\).
The annex proves both smoothness and the obstruction to any
well-founded representation, not merely the presence of one descending
chain in this particular model.
\end{hbtheorem}

\subsubsection{Exception ranks and the closure construction}

\begin{hbdefinition}{def:p08-rational-exception-ranks}{Preferential exception ranks of a knowledge base}
Let \(K\) be a set of assertions and \(K^p\) its preferential
closure. A formula \(\alpha\) is exceptional for \(K\) iff
\(\top\mathrel{|\!\sim}\neg\alpha\in K^p\).
An assertion is exceptional when its antecedent is. Let \(E(K)\)
be the exceptional assertions of \(K\). This operator \(E(K)\)
is not the set \(E\) of formula-equivalence classes above.
Define \(C_0=K\), \(C_{\tau+1}=E(C_\tau)\), and, at a limit,
\(C_\tau=\bigcap_{\rho<\tau}C_\rho\).
The rank of \(\alpha\) for \(K\) is the least ordinal \(\tau\)
at which it is not exceptional for \(C_\tau\); otherwise it has
no rank. A base is admissible when every no-rank formula satisfies
\(\alpha\mathrel{|\!\sim}\bot\in K^p\).
The material counterpart \(\widetilde{\mathcal A}\) of
\(\mathcal A=(\alpha\mathrel{|\!\sim}\beta)\) is
\(\alpha\to\beta\); \(\widetilde K\) is the set of all these
counterparts. This is an operation on assertions, not identification
of material implication with defeasible consequence.
\end{hbdefinition}

\begin{hbdefinition}{def:p08-rational-abstract-closure}{The source's selection criterion for rational closure}
For a rational relation \(R\), the assertion abbreviated
\(\alpha<\beta\) in the source's closure comparison is
\(\alpha\vee\beta\mathrel{|\!\sim}\neg\beta\in R\).
This abbreviation is distinct from the class order used above; in
particular inconsistent formulas require care. For rational relations
\(K_0,K_1\), write \(K_0\prec K_1\) when both conditions hold:
there is \(\alpha\mathrel{|\!\sim}\beta\in K_1\setminus K_0\)
such that every assertion of \(K_0\) whose antecedent \(\gamma\)
satisfies \(\gamma<\alpha\) in \(K_0\) also lies in \(K_1\);
and, for each \(\gamma\mathrel{|\!\sim}\delta\in K_0\setminus K_1\),
there is \(\rho\mathrel{|\!\sim}\eta\in K_1\setminus K_0\)
with \(\rho<\gamma\) in \(K_1\).
The rational closure \(\overline K\), when it exists, is a rational
extension of \(K\) preferable in this sense to every other rational
extension. It is not the intersection of all rational extensions or
consequence in all ranked models of \(K\).
\end{hbdefinition}

\begin{hbtheorem}{thm:p08-rational-source-selection}{proved=annex:p08.rational.selection}{The admissible-base selection theorem}
For an admissible base, the source's rational closure exists uniquely and consists
exactly of assertions \(\alpha\mathrel{|\!\sim}\beta\) for
which \(\alpha\) has no rank, or its rank is strictly smaller than
that of \(\alpha\wedge\neg\beta\), including the case in which
the latter has no rank. This is source Theorem~7, pp.~34--36; the
strict ordering of rational relations is source Lemma~25, pp.~29--31.
The annex proves that the criterion is rational, extends the base, and
is preferable to each other rational extension under both quantified
comparison conditions. The finite ranked realization and algorithm
below have separate full proofs. Arbitrary infinite bases need not be
admissible and need not
have a rational closure; source Example~1 and Lemma~27 concern that
separate obstruction.
\end{hbtheorem}

\begin{hbtheorem}{thm:p08-rational-material-exception}{proved=annex:p08.rational.material}{The exact classical exception test, for arbitrary bases}
For any set-sized \(K\) and formula \(\delta\),
\(\widetilde K\models\delta\) iff
\(\top\mathrel{|\!\sim}\delta\in K^p\).
Thus \(\alpha\) is exceptional iff
\(\widetilde K\models\neg\alpha\).
One implication uses classical compactness and the derived S rule;
the other gives an explicit one-state countermodel. The result does
not say that all preferential entailment reduces directly to the same
material entailment test.
\end{hbtheorem}

\begin{hbtheorem}{thm:p08-rational-finite-admissibility}{proved=annex:p08.rational.admissibility}{Finite bases are admissible, including persistent exceptions}
If \(K\) has \(n\) assertions, its decreasing sequence stabilizes
at some \(C_N\), \(N\leq n\). A formula has no rank iff it is
exceptional for this stable base. Every smooth preferential model of
\(C_N\) contains no state satisfying an antecedent in \(C_N\),
and no state satisfying a no-rank formula. Consequently every no-rank
formula is inconsistent for \(K^p\), and \(K\) is admissible.
The proof permits \(C_N\ne\varnothing\); it does not discard an
inconsistent cycle of defaults or treat it as a ranked world of
``infinite height.''
\end{hbtheorem}

\begin{hbdefinition}{def:p08-rational-finite-realization}{A concrete ranking for a finite base}
For finite \(K\), write \(h(\alpha)\) for its exception rank,
using a formal \(\infty\) when it has no rank. Order this marker
above every finite rank; it is not an extra state rank.
Let
\[
 \mathcal U_* =\{m\in\mathcal U:m\models\widetilde C_N\},\qquad
 k(m)=\min\{i\leq N:m\models\widetilde C_i\}\quad(m\in\mathcal U_*).
\]
Define \(W_K\) with one state for each \(m\in\mathcal U_*\),
identity labeling, and preference \(m\prec m'\) iff \(k(m)<k(m')\).
The finite number of layers guarantees least ranks for every nonempty
formula set even when \(\mathcal U\) is infinite. A missing valuation
is not silently retained at rank \(\infty\).
\end{hbdefinition}

\calculussheet{sheet:p08-rational-closure}{Lehmann--Magidor rational relations and finite rational-closure procedure}{Flat propositional assertions $\alpha\mathrel{|\!\sim}\beta$ over the full valuation universe $\mathcal U$; $\to$ is material implication; classical negation is internal to formulas; no nesting of assertions or equality predicate; the six P rules plus Rational Monotonicity characterize a whole relation, while the finite closure procedure uses the separately defined exception operator $E(C)$.}
\hbrule{rule:p08-rational-lle}{Left Logical Equivalence}{\models\alpha\leftrightarrow\beta\qquad\alpha\mathrel{|\!\sim}\gamma}{\beta\mathrel{|\!\sim}\gamma}{classical validity over all $\mathcal U$}
\hbrule{rule:p08-rational-rw}{Right Weakening}{\models\beta\to\gamma\qquad\alpha\mathrel{|\!\sim}\beta}{\alpha\mathrel{|\!\sim}\gamma}{classical validity over all $\mathcal U$}
\hbrule{rule:p08-rational-ref}{Reflexivity}{}{\alpha\mathrel{|\!\sim}\alpha}{\none}
\hbrule{rule:p08-rational-and}{And}{\alpha\mathrel{|\!\sim}\beta\qquad\alpha\mathrel{|\!\sim}\gamma}{\alpha\mathrel{|\!\sim}\beta\wedge\gamma}{\none}
\hbrule{rule:p08-rational-or}{Or}{\alpha\mathrel{|\!\sim}\gamma\qquad\beta\mathrel{|\!\sim}\gamma}{\alpha\vee\beta\mathrel{|\!\sim}\gamma}{\none}
\hbrule{rule:p08-rational-cm}{Cautious Monotonicity}{\alpha\mathrel{|\!\sim}\beta\qquad\alpha\mathrel{|\!\sim}\gamma}{\alpha\wedge\beta\mathrel{|\!\sim}\gamma}{\none}
\hbrule{rule:p08-rational-rm}{Rational Monotonicity}{\alpha\mathrel{|\!\sim}\gamma\qquad\alpha\mathrel{|\!\not\sim}\neg\beta}{\alpha\wedge\beta\mathrel{|\!\sim}\gamma}{the negative premise is actual nonmembership in the complete relation, not failed finite proof search; soundness is Annex~\ref{annex:p08.rational.soundness}}
\hbrule{rule:p08-rational-s}{S}{\alpha\wedge\beta\mathrel{|\!\sim}\gamma}{\alpha\mathrel{|\!\sim}\beta\to\gamma}{derived P rule; Annex~\ref{annex:p08.klm.1}}
\hbrule{rule:p08-rational-cut}{Cut}{\alpha\mathrel{|\!\sim}\beta\qquad\alpha\wedge\beta\mathrel{|\!\sim}\gamma}{\alpha\mathrel{|\!\sim}\gamma}{derived P rule; Annex~\ref{annex:p08.klm.1}; this is not unrestricted transitivity or arbitrary left strengthening}
\hbrule{rule:p08-rational-14}{Source (14)}{\alpha\mathrel{|\!\sim}\bot}{\alpha\wedge\beta\mathrel{|\!\sim}\bot}{derived P rule; Annex~\ref{annex:p08.rational.transport}}
\hbrule{rule:p08-rational-15}{Source (15)}{\alpha\vee\beta\mathrel{|\!\sim}\neg\beta}{\alpha\mathrel{|\!\sim}\neg\beta}{derived P rule; Annex~\ref{annex:p08.rational.transport}}
\hbrule{rule:p08-rational-16}{Source (16)}{\alpha\vee\beta\vee\gamma\mathrel{|\!\sim}\neg\alpha\wedge\neg\beta}{\beta\vee\gamma\mathrel{|\!\sim}\neg\beta}{derived P rule; Annex~\ref{annex:p08.rational.transport}}
\hbrule{rule:p08-rational-17}{Source (17)}{\alpha\vee\beta\mathrel{|\!\sim}\neg\alpha}{\alpha\vee\beta\vee\gamma\mathrel{|\!\sim}\neg\alpha}{derived P rule; Annex~\ref{annex:p08.rational.transport}}
\hbrule{rule:p08-rational-18}{Source (18)}{\alpha\vee\gamma\mathrel{|\!\sim}\neg\alpha\qquad\beta\vee\gamma\mathrel{|\!\not\sim}\neg\beta}{\alpha\vee\beta\mathrel{|\!\sim}\neg\alpha}{derived rational comparison rule; Annex~\ref{annex:p08.rational.transport}}
\hbrule{rule:p08-rational-material}{Exception test}{\widetilde C\models\neg\delta}{\text{$\delta$ is exceptional for $C$}}{equivalence in Annex~\ref{annex:p08.rational.material}; test unsatisfiability of $\widetilde C\cup\{\delta\}$}
\hbrule{rule:p08-rational-start}{Procedure: initialize}{}{C:=K}{finite $K$ and query $\alpha\mathrel{|\!\sim}\beta$}
\hbrule{rule:p08-rational-next}{Procedure: advance}{\widetilde C\models\neg\alpha\qquad E(C)\ne C}{C:=E(C)}{repeat while both guards hold; this is a state update, not an additional inference rule of R}
\hbrule{rule:p08-rational-yes}{Procedure: yes}{\widetilde C\models\neg(\alpha\wedge\neg\beta)}{\text{answer yes}}{after the loop has stopped}
\hbrule{rule:p08-rational-no}{Procedure: no}{\widetilde C\not\models\neg(\alpha\wedge\neg\beta)}{\text{answer no}}{after the loop has stopped; full termination and correctness in Annex~\ref{annex:p08.rational.algorithm}}


\begin{hbtheorem}{thm:p08-rational-finite-realization}{proved=annex:p08.rational.realization}{Rank criterion, rationality and exact realization}
For each formula, \(h(\alpha)\) is the least \(k\)-rank of its
states in \(W_K\), or \(\infty\) exactly when it has no states.
The model satisfies \(K\), is smooth and ranked, and represents
precisely the relation
\[
 \alpha\mathrel{|\!\sim}\beta
 \quad\Longleftrightarrow\quad
 h(\alpha)=\infty\ \text{ or }
 h(\alpha)<h(\alpha\wedge\neg\beta).
\]
In particular this relation is rational. The source selection theorem
identifies it with \(\overline K\); the realization proof separately
establishes every equality and model condition in this display.
\end{hbtheorem}

\begin{hbtheorem}{thm:p08-rational-finite-algorithm}{proved=annex:p08.rational.algorithm}{Termination and correctness of the finite query procedure}
Starting with \(C=K\), while \(\alpha\) is exceptional for \(C\)
and \(E(C)\ne C\), replace \(C\) by \(E(C)\).
Return yes iff \(\alpha\wedge\neg\beta\) is exceptional for
the resulting \(C\). This decides the displayed rank criterion,
including a no-rank antecedent. Each exception test is classical
unsatisfiability of \(\widetilde C\cup\{\delta\}\) for the
tested formula \(\delta\). At most \(O(n^2)\) such tests are
needed for \(n\) defaults. This is a polynomial number of SAT
queries, not a deterministic polynomial-time bound for unrestricted
propositional input. If every actual tested conjunction is Horn, Horn
satisfiability gives a polynomial procedure; merely calling the
defaults ``Horn'' without checking the query tests does not suffice.
\end{hbtheorem}

\begin{hbtheorem}{thm:p08-rational-ranked-intersection}{proved=annex:p08.rational.intersection}{Why all-ranked-model consequence is not rational closure}
For \(K=\{p\mathrel{|\!\sim}q\}\) on atoms \(p,q,a\),
rational closure includes \(p\wedge a\mathrel{|\!\sim}q\).
There is nevertheless a finite smooth ranked model of \(K\)
refuting it. More strongly, the intersection of the relations of two
finite ranked models can fail Rational Monotonicity. Both models
satisfy \(p\mathrel{|\!\sim}q\); their intersection contains
this assertion but not \(p\mathrel{|\!\sim}\neg a\), and
does not contain \(p\wedge a\mathrel{|\!\sim}q\).
Thus all-ranked-model consequence need not be a rational relation.
\end{hbtheorem}

\begin{hbtheorem}{thm:p08-rational-pointwise-minimum}{proved=annex:p08.rational.minimum}{The lowest complete finite valuation ranking}
On any finite vocabulary containing the atoms of \(K\),
\(\mathcal U_*\) is exactly the valuations that can label states
in a smooth preferential model of \(K\). Among ranked models
containing each of these valuations, normalize their finitely many
occupied layers to nonnegative integers, and call their least-copy
valuation ranking \(f(m)\). The rank \(k(m)\) above
is pointwise no greater than the least rank of any copy of \(m\)
in any such model. It is therefore the unique least complete valuation
ranking. This pointwise model statement has a full induction proof;
it is not substituted for the source's different ordering of whole
rational consequence relations.
\end{hbtheorem}

\begin{hbtheorem}{thm:p08-rational-penguin-example}{proved=annex:p08.rational.penguin}{All layers of the penguin triangle}
For \(K=\{p\mathrel{|\!\sim}b,
p\mathrel{|\!\sim}\neg f,b\mathrel{|\!\sim}f\}\),
\(C_1=\{p\mathrel{|\!\sim}b,p\mathrel{|\!\sim}\neg f\}\)
and \(C_2=\varnothing\).
Over \(b,p,f\), the rank-0 valuations have \(p=0\) and
\(b\to f\); rank 1 consists of \((b,p,f)=(1,1,0)\) and
\((1,0,0)\); rank 2 consists of \(p=1\) with
\(\neg(b\wedge\neg f)\). The coordinates are written in
the order \((b,p,f)\).
Thus \(h(b)=0\), \(h(p)=1\),
\(h(p\wedge f)=2\), and both \(p\mathrel{|\!\sim}\neg f\)
and \(p\wedge a\mathrel{|\!\sim}\neg f\) hold for a fresh
atom \(a\), whereas \(p\mathrel{|\!\sim}f\) fails.
Neither \(b\wedge\neg f\mathrel{|\!\sim}p\) nor
\(b\wedge\neg f\mathrel{|\!\sim}\neg p\) holds.
An independent default \(b\mathrel{|\!\sim}a\) is not inherited
by \(p\): the construction then has minimal \(p\)-states of
both values of \(a\). This is the loss of generic-property
inheritance to exceptional subclasses, not failure of rationality.
\end{hbtheorem}

\begin{hbtheorem}{thm:p08-rational-nixon-example}{proved=annex:p08.rational.nixon}{Conflicting defaults and a persistent contradictory base}
For \(K=\{r\mathrel{|\!\sim}\neg q,
s\mathrel{|\!\sim}q\}\), neither default is exceptional,
so \(C_1=\varnothing\). The conjunction \(r\wedge s\)
has rank 1; each of its conjunctions with \(q\) and \(\neg q\)
also has rank 1. Rational closure therefore chooses neither consequent
for \(r\wedge s\). A fresh \(a\) nevertheless permits
\(r\wedge a\mathrel{|\!\sim}\neg q\).
In contrast, for \(K'=\{p\mathrel{|\!\sim}q,
p\mathrel{|\!\sim}\neg q\}\), the base is already stable,
\(p\) has no rank, all realized states satisfy \(\neg p\),
and every assertion with antecedent \(p\) holds vacuously.
The two cases must not be identified as the same unresolved conflict.
\end{hbtheorem}

Rational Monotonicity is a constraint on a complete conditional relation,
not a rule for accepting a conclusion whenever a program has not yet
derived its negation. The arbitrary-language representation above does
not yield an algorithm for arbitrary infinite bases. The finite-base
algorithm does not make rational closure an empirically universal account
of human reasoning, a nested counterfactual logic, an AGM revision
operation, or a probability-threshold rule. C, CL, CM and M retain their
separate cumulative, loop, monotonicity and model-label hypotheses;
none of their representation constructions follows from the ranked
construction \cite{KrausLehmannMagidor2002Author,LehmannMagidor2002Author}.
